Let $H\cup H^\prime$ be a Heegaard splitting of a closed orientable 3-manifold.    Let $\mathcal{C}$ denote the collection of closed curves in $\bd H$ that bound a disc in $H^{\prime}$.  Does the set of algebraically simplest elements of $\mathcal{C}$ overlap the set of geometrically simplest elements?  Below we make this question precise and begin an answer to it.

Let $\gamma$ be a closed curve in a handlebody $H$.  Define the girth of $\gamma$ to be the smallest intersection number between a curve homotopic to $\gamma$ and an essential disc in $H$.  Girth is an algebraic measure of complexity.  We say $\gamma$ is \emph{primitive} if $\gamma$ does not represent a proper power of some element of $\pi_{1}(H)$.  In particular, primitive curves are nontrivial in $\pi_{1}(H)$.  We propose the following answer to the above question.

\begin{girth conjecture}
Let $H\cup H^\prime$ be a Heegaard splitting of a closed orientable 
3-manifold.  Let $\gamma$ be a closed curve in $\bd H$ that is
primitive in $H$ and contracts in $H^\prime$.  Then there is an essential simple closed curve $\gamma^{\prime}$ in $\bd H$ with girth no greater than that of $\gamma$ which contracts in $H^\prime$.
\end{girth conjecture}

For example, let $M$ be a closed hyperbolic manifold.  Let $H\cup H^{\prime}$ be a Heegaard splitting of $M$.  Let $\gamma$ be an immersed curve in $\bd H$ that bounds a disc in $H^{\prime}$.  Does $\bd H$ contain an embedded curve $\gamma^{\prime}$ that bounds a disc in $H^{\prime}$ and has girth no greater than that of $\gamma$?  The Girth Conjecture would say yes.

Commonly in 3-manifold topology we wish to construct embedded objects from immersed objects, and Heegaard splittings are a fundamental setting.  The Girth Conjecture is like a Poincare Conjecture for all 3-manifolds.  By Proposition~\ref{finite fundamental group}, the Girth Conjecture for 3-manifolds with finite fundamental group is equivalent to the spherical space form conjecture.

Splittings of reducible 3-manifolds satisfy the Girth Conjecture by Lemma~\ref{reducible} and work of~\cite{Ha}.  In addition to spherical 3-manifolds, solvmanifolds satisfy the Girth Conjecture~\cite{Je1}, as well as the simplest genus two hyperbolic manifolds~\cite{Je2}.  See Figure~\ref{whitehead}.  The classification of Heegaard splittings of Seifert-fibered spaces~\cite{MS} provides an answer for genus two splittings of Seifert-fibered spaces.  We discuss connections to the geometry of 3-manifolds in more detail in Section~\ref{heegaardgirth}.

In Section~\ref{heegaardgirth}, we define the algebraic and geometric girths of a Heegaard splitting.  We show in Proposition~\ref{algebraic girth one} that the Girth Conjecture holds for all Heegaard splittings of algebraic girth one.  The proof uses the Freiheitssatz and Kneser's Conjecture for 3-manifolds.

Our Main Theorem shows that the Girth Conjecture is true in a certain general setting.  It is a purely topological result.

\begin{figure}
\epsfig{file=figures/whitehead.eps, height=2.5in}
\caption{
All Dehn fillings of this standard Heegaard splitting $H\cup H^\prime$
of the Whitehead link $L$ satisfy 
Conjecture~\ref{ag girth conjecture}.
All hyperbolic manifolds that result have algebraic and
geometric girth four~\cite{Je2}. 
The thin curve bounds a disc in the gray handlebody $H$ that is disjoint from $L$.  This disc visibly intersects a disc in the complement $H^\prime$ 
in four points.}
\label{whitehead}
\end{figure}

\begin{subsurface theorem}[\textbf{Theorem~\ref{subsurface theorem}}]
Let $H$ be a handlebody.  Let $S$ be a compact surface in 
$\partial H$.  The minimum girth in $H$ over all essential closed
curves in $S$ can be achieved by a simple closed curve in $S$. 
\end{subsurface theorem}

The Main Theorem implies the Girth Conjecture for curves $\gamma$ that are contained in an embedded ball in $H^\prime$.  If $B$ is a closed ball in $H^{\prime}$ that contains $\gamma$, simply take $S$ to be $B\cap \bd H$.  All curves in $B\cap \bd H$ contract in $H^{\prime}$.  Theorem~\ref{subsurface theorem} is trivial for surfaces $S$ that are compressible in $H$.  More interesting and typical surfaces $S$ can be obtained by taking the complement of a disjoint collection of simple closed curves in $\bd H$.

Our Main Theorem has the following corollaries.
Let $H$ be a handlebody.
A \emph{disc system} for $H$ is a disjoint collection of essential 
discs in $H$ that divide $H$ into a ball.
A Heegaard splitting $H\cup H^\prime$ is \emph{weakly reducible}
if there are essential discs in $H$ and $H^\prime$ whose
boundaries do not intersect in $\bd H$.  By ~\cite{CG}, weak
reducibility implies that either the 3-manifold is Haken,
or the splitting is reducible.

\begin{corollary}
\label{cor 1}
Let $H\cup H^\prime$ be a Heegaard splitting of a closed orientable 
3-manifold $M$.
Let $\delta$ be a collection of simple closed curves in $\partial H$ 
that bounds a disc system in $H^\prime$.
If there exists a curve in $\partial H-\delta$ that is homotopic in 
$H$ to a curve disjoint from some essential disc in $H$, 
then either $M$ is Haken, or the splitting is reducible.
\end{corollary}

\begin{proof}
Let $\gamma$ be a curve in $\partial H-\delta$ that is
homotopic in $H$ to a curve disjoint from an essential disc in $H$.
Then $\gamma$ has girth zero in $H$.  By the Main Theorem,
$\partial H-\delta$ contains a simple closed curve of girth zero
in $H$.

So we can assume $\gamma$ is simple.  By 
Lemma~\ref{one embedded curve}, the curve $\gamma$ is disjoint from 
an essential disc $D$ in $H$.
Since $\gamma$ is disjoint from a disc system in $H^\prime$,
$\gamma$ bounds an embedded disc $D^\prime$ in $H^\prime$.

Since $D$ and $D^\prime$ are essential discs in $H$ and $H^\prime$
with disjoint boundaries, the splitting $H\cup H^\prime$ 
is weakly reducible.  By~\cite{CG}, either 
the manifold is Haken, or the Heegaard splitting is reducible.
\end{proof}

The distance of a Heegaard splitting is one measure of 
complexity for Heegaard splittings.  Distance is explored,
for instance, in~\cite{He2}.
A distance one splitting is the same as a weakly reducible splitting.
Distance two also has special significance.
A Heegaard splitting $H\cup H^\prime$ has distance two if 
there exists a curve $\gamma$ in $\bd H$ that is disjoint from an 
essential disc in both $H$ and $H^\prime$.
This is also called the 
\emph{disjoint curve property}~\cite{Tho}.
 
\begin{corollary}
Let $H\cup H^\prime$ be a Heegaard splitting of a closed orientable 
3-manifold.
Let $D^\prime$ be an essential disc in $H^\prime$.
If there is a curve in $\partial H-\bd D^\prime$ that is homotopic 
in $H$ to a curve disjoint from some essential disc in $H$, then 
$H\cup H^\prime$ satisfies the disjoint curve property.
\end{corollary}

\begin{proof}
Let $\gamma$ be a curve in $\partial H-\bd D^\prime$ that is
homotopic in $H$ to a curve disjoint from some essential disc.
Then $\gamma$ has girth zero in $H$.  By the Main Theorem,
$\bd H-\bd D^\prime$ contains a simple closed curve of girth zero.

So assume that $\gamma$ is simple.
By Lemma~\ref{one embedded curve}, the curve $\gamma$ is disjoint 
from some essential disc $D$ in $H$.
Since the curve $\gamma$ is disjoint from both $D$ and $D^\prime$,
the splitting $H\cup H^\prime$ has distance two.
\end{proof}

We suggest two approaches to the Girth Conjecture for further
study.  Consider a Heegaard splitting $H\cup H^\prime$.
Let $\gamma$ be an immersed curve in $\bd H$ 
that bounds a disc $D^\prime$ in $H^\prime$.  Can we make 
$\gamma$ simple without increasing its girth in $H$?
The tower construction is one method of producing an 
embedded disc from an immersed curve.  Papakyriakapolous
used this method to prove the Loop Theorem.
See~\cite[p.45]{Hat} or~\cite[p.7]{Ja} for a description
of this construction.

If a neighborhood of the disc $D^\prime$ is a ball, we can 
apply the Main Theorem to $\gamma$.  If not, there 
are nontrivial coverings of this neighborhood.  The tower construction
proceeds by factoring the disc map through 
a sequence of 2-sheeted coverings.
This inspires the following question.

\begin{question}
Let $H\cup H^\prime$ be a Heegaard splitting of a closed $3$-manifold.
Let $S$ be a surface immersed in $\partial H^\prime$ so that $\pi_1(S)$
maps trivially into $H^\prime$.  Consider the minimum girth in $H$ 
over all essential closed curves in $S$.  Can the minimum be 
achieved by a simple closed curve in $S$?
\end{question} 

For embedded $S$ the answer is yes by the Main Theorem.
For $S$ an annulus the answer is yes by Lemma~\ref{power}.  
The answer is no if $S$ does not necessarily contract in $H^{\prime}$.  
To answer this question it suffices to consider the case of $S$
a thrice-punctured sphere.

Here is a second approach.
Let $S$ be a compact surface with boundary $\bd S$.
Let $\nn(\bd S)$ denote the subgroup of $\pi_1(S)$ normally 
generated by the boundary curves in $\bd S$.  This 
parallels the definition of the normal subgroup of relators 
in a Heegaard diagram.
The following statement resembles both the Girth Conjecture 
and the Main Theorem.

\begin{conjecture}
Let $H$ be a handlebody.  Let $S$ be a compact surface in $\bd H$.
Then the minimum girth in $H$ over all closed curves 
in $\nn(\bd S)$ can be achieved by a simple closed curve in $\nn(\bd S)$.
\end{conjecture}

In Section~\ref{trees} we introduce some basic results and constructions related to trees and handlebodies, and in Section~\ref{free groups} we do the same for free groups.  In Section~\ref{handlebodygirth} we present some elementary properties of girth.  We focus on simple closed curves.  There we also introduce a special conjecture, Conjecture~\ref{special conjecture}. In Section~\ref{heegaardgirth} we discuss girth for Heegaard splittings.  We define the algebraic and geometric girths of 
a Heegaard splitting and present a statement equivalent to the 
Girth Conjecture.  We also describe some results relating 
the girth of Heegaard splittings to 3-manifold geometry.

In Section~\ref{subordination} we prove Lemma~\ref{subordination-lemma}, an important lemma that we use several times throughout the paper.
In Section~\ref{curvealgorithm} we present an algorithm 
to compute the girth of a curve.  The algorithm uses Lemma~\ref{subordination-lemma}.

In Sections~\ref{maintheorem} and~\ref{mainlemma} we prove 
the Main Theorem.  We present several lemmas and the 
main proof in Section~\ref{maintheorem}.  We reserve the 
key Lemma~\ref{main lemma} for Section~\ref{mainlemma}.
The proof of Lemma~\ref{main lemma} contains the central construction 
used in the proof of the Subsurface Theorem.
Definitions needed to understand the proofs are spread throughout 
Sections~\ref{subordination} through~\ref{mainlemma}.

We conclude with an algorithm in Section~\ref{surfacealgorithm} to compute the girth for surfaces.  This short section pulls together results from several sections of the paper.